% R1082. Tríptico de gravuras originais; as fontes raster são preservadas.
\def\Pages{436}
\def\Spine{25.99161680}
\def\CoverWidth{352.34161680}

\documentclass{article}
\usepackage[paperwidth=\CoverWidth mm,paperheight=236.35mm,margin=0pt]{geometry}
\usepackage{fontspec,graphicx,tikz,xcolor}
\setmainfont{EB Garamond}
\setsansfont{Lato}
\usepackage[pdfencoding=auto,unicode]{hyperref}
\hypersetup{pdftitle={Molambudos: Piranesi e Pencz, capa \Pages\ páginas},pdfauthor={Marcelo Dias de Carvalho Filho},pdfsubject={Brochura 160 x 230 mm; cor premium em papel branco; sangria 3,175 mm}}
\pagestyle{empty}
% Swatches CMYK nominais <= 210%; não constitui perfil ICC ou TAC global.
\definecolor{paper}{cmyk}{0,0.04132,0.14463,0.05098}
\definecolor{ink}{cmyk}{0,0.23256,0.46512,0.83137}
\definecolor{ruleink}{cmyk}{0,0.18,0.35,0.52}
\begin{document}
\null
\begin{tikzpicture}[remember picture,overlay]
\begin{scope}[shift={(current page.north west)},x=1mm,y=-1mm]
\pgfmathsetmacro{\FrontX}{163.175+\Spine}
\pgfmathsetmacro{\FrontCenter}{\FrontX+80}
\pgfmathsetmacro{\SpineCenter}{163.175+\Spine/2}
\fill[paper] (0,0) rectangle (\CoverWidth,236.35);

% CONTRACAPA
\draw[ruleink,line width=.75pt] (19.175,19) -- (147.175,19);
\node[anchor=north west,inner sep=0,text=ink,align=left,font=\fontsize{20}{22}\selectfont] at (19.175,28) {HÁ HISTÓRIAS QUE\\O TEMPO ENTERRA.\\MAS NÃO APAGA.};
\node[anchor=north west,inner sep=0,text=ink,text width=126mm,align=left,font=\fontsize{12.5}{16}\selectfont] at (19.175,63) {\hyphenpenalty=10000\exhyphenpenalty=10000\raggedright
Em 1915, a seca arranca Joaquim Antônio Correia de seu sertão. Anos depois, reduzido ao número 1.260 em uma instituição psiquiátrica, ele escreve para não desaparecer. Suas memórias atravessam a fome, o confinamento e os silêncios que nenhum prontuário consegue explicar.

\vspace{3mm}
Décadas mais tarde, a psiquiatra forense Lúcia Mendes reúne diários, laudos, cartas e fotografias para reconstituir uma história que parece resistir ao esquecimento. Mas cada fragmento abre outro caminho, e o arquivo começa a interpelar quem o lê.

\vspace{3mm}
Em \textit{Molambudos}, o horror nasce da memória ferida — e continua onde termina a página.
};
\node[anchor=north west,inner sep=0,text=ink,align=left,font=\fontsize{15}{18}\selectfont\itshape] at (19.175,145) {O ciclo recomeça a cada leitor.};
\node[anchor=north west,inner sep=0] at (19.175,165) {\includegraphics[width=42mm]{arte_arquivo.png}};
\node[anchor=north west,inner sep=0,text=ink,align=left,font=\sffamily\fontsize{8.5}{12}\selectfont] at (78,174) {ROMANCE BRASILEIRO\\HORROR PSICOLÓGICO};
% Reserva branca e vazia para código automático, 50,8 x 30,48 mm.
\fill[white] (106.025,196.345) rectangle (156.825,226.825);

% LOMBADA
\draw[ruleink,line width=.75pt] (\SpineCenter,18) -- (\SpineCenter,30);
\node[rotate=-90,inner sep=0,text=ink,font=\fontsize{22}{24}\selectfont\bfseries] at (\SpineCenter,82) {MOLAMBUDOS};
\node[rotate=-90,inner sep=0,text=ink,font=\fontsize{9.5}{12}\selectfont] at (\SpineCenter,158) {MARCELO DIAS DE CARVALHO FILHO};
\node[inner sep=0,text=ink,align=center,font=\sffamily\fontsize{7.3}{9}\selectfont] at (\SpineCenter,213) {OPENCODE\\BOOKS};

% FRENTE: três gravuras de 42 x 63 mm; 1024 x 1536 px = 619 dpi nativos.
\draw[ruleink,line width=.75pt] (\FrontX+14,17) -- (\FrontX+146,17);
\node[inner sep=0,text=ink] at (\FrontCenter,33) {\resizebox{136mm}{!}{\bfseries MOLAMBUDOS}};
\node[inner sep=0,text=ink,font=\fontsize{16}{19}\selectfont] at (\FrontCenter,50) {O Diário do Paciente 1.260};
\node[inner sep=0,text=ink,font=\sffamily\addfontfeatures{LetterSpace=14}\fontsize{7.8}{10}\selectfont] at (\FrontCenter,62) {UM ROMANCE EM FRAGMENTOS};
\node[anchor=north west,inner sep=0] at (\FrontX+13,78) {\includegraphics[width=42mm]{arte_sertao.png}};
\node[anchor=north west,inner sep=0] at (\FrontX+59,78) {\includegraphics[width=42mm]{arte_confinamento.png}};
\node[anchor=north west,inner sep=0] at (\FrontX+105,78) {\includegraphics[width=42mm]{arte_arquivo.png}};
\node[inner sep=0,text=ink,font=\scshape\fontsize{8.5}{11}\selectfont] at (\FrontX+34,149) {Memória};
\node[inner sep=0,text=ink,font=\scshape\fontsize{8.5}{11}\selectfont] at (\FrontX+80,149) {Confinamento};
\node[inner sep=0,text=ink,font=\scshape\fontsize{8.5}{11}\selectfont] at (\FrontX+126,149) {Arquivo};
\node[inner sep=0,text=ruleink,font=\sffamily\fontsize{7.3}{10}\selectfont] at (\FrontCenter,160) {RECONSTITUIÇÕES FICCIONAIS};
\node[inner sep=0,text=ink,align=center,font=\fontsize{12.5}{15}\selectfont\itshape] at (\FrontCenter,180) {Ela não é um monstro.\\Ela é uma consequência.};
\node[inner sep=0,text=ink,font=\fontsize{13.4}{16}\selectfont] at (\FrontCenter,201) {Marcelo Dias de Carvalho Filho};
\node[inner sep=0,text=ink,font=\fontsize{10}{13}\selectfont\itshape] at (\FrontCenter,210) {Coautor e ilustrador: Marcelo Claro};
\draw[ruleink,line width=.75pt] (\FrontX+14,217.5) -- (\FrontX+146,217.5);
\node[inner sep=0,text=ink,font=\sffamily\addfontfeatures{LetterSpace=12}\fontsize{7.2}{9}\selectfont] at (\FrontCenter,223) {OPENCODE BOOKS};
\end{scope}
\end{tikzpicture}
\end{document}
